% =============================================================================================
% SECTION 4 -- CONCLUSIONS.  Draft 2.
%
% Follows the blueprint's conclusion shape: restate the stakes, restate the bottleneck, name the
% need, say what we did, say what validates it, say what it is for, and save the payoff line for
% last. Their conclusion carries counts but NO error statistics -- the numbers live in the
% abstract and the results table, and repeating them here reads as a summary rather than a close.
% Ours keeps two (the offset and the within-2x rate) because they are the claim itself, and drops
% the other four the previous draft restated.
%
% The last sentence is the one thing here that appears nowhere else in the paper.
% =============================================================================================

\section{Conclusions}
\label{sec:conclusions}

Lattice thermal conductivity is the quantity that decides whether a half Heusler is worth
synthesising, and it is also the quantity first-principles transport theory is most often asked to
supply. Because converged calculations are expensive, the field has increasingly reproduced them
with machine-learned surrogates trained on calculated labels---which are what exist in
quantity. The result is a literature of models that predict a calculation accurately and a
measurement only incidentally.

We have shown that the two scales differ by a knowable amount. Across half Heuslers for which both a
published transport calculation and a laboratory measurement exist, the calculation exceeds the
measurement by a median factor of \num{1.55}, one-sidedly and in every bonding family, and the
offset is regular enough to be removed by a closed-form correction with two parameters. Fitting that
correction required a corpus that does not average the two quantities together, so we built one:
\num{52} measured half Heuslers drawn from \num{164} publications alongside \num{289} carrying
transport calculations, every value tiered by the method that produced it.

Validated with each compound's entire chemistry withheld from the model and from the calibration
alike, and inside a domain of applicability declared before the validation set was scored, the
correction places \SI{86.8}{\percent} of held-out compounds within a factor of two of measurement
and orders them reliably. Outside that domain it fails, by a factor we report rather than conceal.
It is a screen and not a curve predictor: it says what a laboratory is likely to measure at a stated
temperature, and its parameters should not be read as a measurement of microstructure. On that basis
we issue five predictions and decline eleven candidates, naming the condition each fails.

The five are falsifiable by a single synthesis each, and we would regard a failure among them as
more informative than the successes. The wider result is that the barrier to predicting more of them
is no longer computational. Fourteen further candidates already carry sound calculations and wait
only on five laboratory measurements across four bonding families---the point at which a
calibration of this kind stops being limited by what has been calculated and starts being limited
by what has been measured.
